\section{What is object-oriented programming?}
\subsection{Procedural programming and its limitations:}
\begin{itemize}
    \item Procedural programming is pretty much what we've been doing throughout the course.
    \item Focus is on processes or actions that a program takes.
    \item Programs are typically a collection of functions.
    \item Data is declared separately.
    \item Data is passed as arguments to into functions.
    \item Fairly easy to learn.
\end{itemize}
Functions need to know the structure of the data.
\begin{itemize}
    \item If the structure of the data changes many functions must be changed.
    \item Imagine a program with thousands of functions that all depend on a data structure being passed in to it, imagine that that data structure needs to change you would need to change every function, functions might even need to be modified to handle such change and a grave impact is in order, then test all the functions which makes development much more expensive.
\end{itemize}

As programs get larger they become more:
\begin{itemize}
    \item Difficult to understand.
    \item Difficult to maintain.
    \item Difficult to extend.
    \item Difficult to debug.
    \item Difficult to reuse code.
    \item Fragile and easier to break.
\end{itemize}

\subsection{What is objected-oriented programming:}
\begin{itemize}
    \item Object oriented denotes many things such as: OO analysis, OO design, OO testing, etcetera.
    \item Object-oriented is all about modeling you software in terms of objects and classes.
    \item Classes and objects:
        \begin{itemize}
            \item Focus is on classes that model real-world domain entities.
            \item Allows developers to think at a higher level of abstraction.
            \item Used successfully in very large programs.
        \end{itemize}
\end{itemize}

Why is it such a big deal?
\begin{itemize}
    \item As our programs grow more and more complex we need ways to deal with complexity, classes and objects are one way to do just that, rather than thinking of details we can think of just that object.
\end{itemize}

Encapsulation:
\begin{itemize}
    \item Objects contain data AND operations that work on that data.
    \item Abstract Data Type (ADT).
\end{itemize}

Information hiding:
\begin{itemize}
    \item OO allows us to hide implementation specific logic in a class that is available only within the class, now we know that the user can mess with the user implementation specific code since they never knew about it in the first place.
    \item Users of the class code to the interface since they don't need to know the implementation.
    \item More abstraction.
    \item Easier to test, debug, maintain and extend.
\end{itemize}

Reusability:
\begin{itemize}
    \item Easier to reuse in other applications.
    \item Faster development.
    \item Higher quality.
\end{itemize}

Inheritance:
\begin{itemize}
    \item Can create new classes in terms of existing classes.
    \item Reusability.
    \item Polymorphic classes.
\end{itemize}

\subsection{OO limitations}
\begin{itemize}
    \item Not a panacea:
        \begin{itemize}
            \item OO programming won't make bad code better.
            \item Not suitable for all types of problems.
            \item Not everything decomposes to a class.
        \end{itemize}
    
    \item Learning curve:
        \begin{itemize}
            \item Usually a steeper learning curve, especially in C++.
            \item Many OO languages, many variations of OO concepts.
        \end{itemize}
    
    \item Design:
        \begin{itemize}
            \item Usually more up-front design is necessary to create good models and hierarchies.
        \end{itemize}
    
    \item Programs can be:
        \begin{itemize}
            \item Larger in size.
            \item Slower.
            \item More complex.
        \end{itemize}
\end{itemize}


%----------------------------------------------------------------------------------------
\section{What are classes and objects?}
\subsection{Classes}
\begin{itemize}
    \item Blueprints from which objects are created.
    \item A user-defined data-type.
    \item Has attributes (data).
    \item Has methods (functions).
    \item Can hide data and methods.
    \item Provides a public interfase.
\end{itemize}
Example classes:
\begin{itemize}
    \item Account.
    \item Employee.
    \item Image.
    \item \mintinline{cpp}{std::vector}.
    \item \mintinline{cpp}{std::string}.
\end{itemize}

\subsection{Objects}
\begin{itemize}
    \item Objects are created from classes.
    \item Represent a specific instantce of a class.
    \item Can create many, many objects.
    \item Each has its own identity.
    \item Each can use the defined class methods. 
\end{itemize}
Examples of objects:
\begin{itemize}
    \item Frank's account is an instance of an account.
    \item Jim's account is an instance of an account.
    \item Each has its own balance, can make deposits, withdrawls, etc.
\end{itemize}

\subsection{Example}
\begin{itemize}
    \item int are not classes but an object is declared in very much the same way any variable is.
\end{itemize}
\begin{minted}[autogobble]{cpp}
    int high_score;
    int low_score;
    Account frank_account; // user defined type.
    Account jim_account;
    std::vector<int> scores;
    std::string name;
\end{minted}


%----------------------------------------------------------------------------------------
\section{Declaring a class and creating objects}
\begin{minted}[autogobble]{cpp}
    #include <iostream>
    class Player {
        // attributes
        std::string name;
        int health;
        int xp;

        // methods
        void talk(std::string text_to_say);
        bool is_dead();
    };
    int main() {
        // Now that we have declared the class, we can create objects or instances of that class.
        Player *enemy = new Player();
        delete enemy;
        return 0;
    }
\end{minted}

Take the following example:
\begin{minted}[autogobble]{cpp}
    #include <iostream>
    class Account {
        // attributes
        std::string name;
        double balance;
        // methods
        bool withdrawl(double amount);
        bool deposit(double amount);
    };
    int main() {
        Account frank_account;
        Account jim_account;
        Account *mary_account = new Account();
        delete mary_account;
        return 0;
    }
\end{minted}

\subsection{Creating objects}
\begin{itemize}
    \item Once we have objects we can use it like any other variable, for example.
        \begin{minted}[autogobble]{cpp}
            Account frank_account, jim_account;
            std::vector<Account> account1 {frank_account, jim_account};
            account1.push_back(jim_account);
            Account accounts[] {frank_account};
        \end{minted}
\end{itemize}

\subsection{Example}
\begin{minted}[autogobble]{cpp}
    #include <iostream>
    #include <string>
    #include <vector>
    using namespace std;
    class Player { // everything inside here is 'encapsulated'.
        // attributes 
        string name;
        unsigned int health;
        int xp;
        // methods prototypes
        void talk(string);
        bool is_dead();
    };
    class Account {
        string name;
        double balance;
        
        bool deposit(double);
        bool withdraw(double);
    };
    int main() {
        Player frank, hero;
        Player * enemy {nullptr};
        enemy = new Player;
        delete enemy;

        Player players[] {frank, hero};

        vector<Player> player_vec;
        player_vec.push_back(hero);

        Account frank_account, jim_account;
        return 0;
    }
\end{minted}

%----------------------------------------------------------------------------------------
\section{Accessing class members}
\begin{itemize}
    \item We can access:
        \begin{itemize}
            \item Class attributes.
            \item Class methods.
        \end{itemize}
    
    \item Some class members will not be accessible (because of the access modifiers).
    \item We need an object to access instance variables.
\end{itemize}

\subsection{Access members}
\begin{itemize}
    \item If we have an object then we can use the dot operator to access the members.
        \begin{minted}[autogobble]{cpp}
            frank_account.balance;
            frank_account.deposit(1000.00);
        \end{minted}
    
    \item If we have a pointer to an object then we can use the \mintinline{cpp}{->} operator or dereference and use the dot operator to access the members.
        \begin{minted}[autogobble]{cpp}
            Account *frank_account = new Account();
        (*frank_account).balance;
        (*frank_account).deposit(1000.00);
        \end{minted}
        \begin{minted}[autogobble]{cpp}
            Account *frank_account = new Account();
            frank_account->balance;
            frank_account->deposit(1000.00);
        \end{minted}
\end{itemize}

\subsection{Example}
\begin{minted}[autogobble]{cpp}
    #include <iostream>
    using namespace std;

    class Player {
        public: // allows us to see the attributes and methods from main.
        string name;
        double health;
        unsigned int xp;

        void talk(string text_to_say) {
            cout << name << " says " << text_to_say << endl;
        }
        bool is_dead(); // not implemented, if executed a linker error is going to be thrown.
    };

    int main() {
        Player frank;
        frank.name = "Frank"; // access the name and change it to "Frank"
        frank.health = 100; // acces the health and change it to 100.
        frank.xp = 12; // access the xp and change it 12.
        frank.talk("Hello");

        Player *enemy = new Player;
        (*enemy).name = "Enemy";
        enemy->name = "Enemy";
        (*enemy).health = 100;
        enemy->health = 100;
        (*enemy).xp = 12;
        enemy->xp = 12;
        (*enemy).talk("I will destroy you.");
        enemy->talk("I will destroy you.");
        
        return 0;
    }
    /* OUTPUT:
    Frank says Hello
    Enemy says I will destroy you.
    Enemy says I will destroy you.
    */
\end{minted}


%----------------------------------------------------------------------------------------
\section{Public and private}
\begin{itemize}
    \item \mintinline{cpp}{public:}
        \begin{itemize}
            \item Accessible everywhere.
        \end{itemize}
    
    \item \mintinline{cpp}{private:}
        \begin{itemize}
            \item Accessible only by memebers or friends (we don't know yet) of the class.
        \end{itemize}
    
    \item \mintinline{cpp}{protected:}
        \begin{itemize}
            \item Used with inheritance (we don't know that yet.)
        \end{itemize}
\end{itemize}

\subsection{Declaration}
\begin{itemize}
    \item Once you declare an access modifier, everything from that point forward will have tha access modifier applied to it.
        \begin{minted}[autogobble]{cpp}
            class Class_name {
                public: 
                // declaration(s);
            };
        \end{minted}
    
    \item By default everything in a class will be private if no access modifier is specified.
        \begin{minted}[autogobble]{cpp}
            class Class_name {
                private: 
                // declaration(s);
            };
        \end{minted}
    
    \item Protected comes in when using inheritance, its basically the private access modifier for inheritance.
        \begin{minted}[autogobble]{cpp}
            class Class_name {
                protected:
                // declaration(s);
            };
        \end{minted}
    
    \item It is common to combine access modifiers because we want to hide only certain things.
        \begin{minted}[autogobble]{cpp}
            class Player {
                private:
                std::string name;
                unsigned int health;
                int xp;
                public:
                void talk(std::string text_to_say);
                bool is_dead();
            };
        \end{minted}
    
    \item If you try to access members declared private you will get a compiler error.
        \begin{minted}[autogobble]{cpp}
            #include <iostream>
            using namespace std;

            class Player {
                private:
                std::string name;
                unsigned int health;
                int xp;
                public:
                void talk(std::string text_to_say);
                bool is_dead();
            };


            int main() {
                Player frank;
                frank.name = "Frank"; // compiler error.
                frank.health = 100; // compiler error.
                frank.talk("Ready to battle.");
                return 0;
            }
        \end{minted}
\end{itemize}


%----------------------------------------------------------------------------------------
\section{Implementing member methods}
\begin{itemize}
    \item Very similar to how we implemented functions.
    \item Member methods have access to the member attributes.
        \begin{itemize}
            \item So you don't need to pass them as arguments.
        \end{itemize}

    \item Can be implemented outside the class declarationg.
        \begin{itemize}
            \item Need to use \verb|class_name::method_name|.
        \end{itemize}
    
    \item Can separate specification from implementation.
        \begin{itemize}
            \item .h file for the class declaration.
            \item .cpp file for the class implementation.
        \end{itemize}
\end{itemize}

\subsection{Inside the class declaration}
\begin{minted}[autogobble]{cpp}
    class Account {
        private:
        double balance;
        public:
        void set_balance(double bal) {
            balance = bal;
        }
        double get_balance() {
            return balance;
        }
    };
\end{minted}

\subsection{Outside the class declaration}
Just have the method prototypes inside the class.
\begin{minted}[autogobble]{cpp}
    class Account {
        private:
        double balance;
        public:
        void set_balance(double bal);
        double get_balance();
    };

    void Account::set_balance(double bal) {
        balance = bal;
    }

    void Account::get_balance() {
        return balance;
    }
\end{minted}

\subsection{Separate specification from implementation}
An include-guard is included so that if you include the file more than once you won't have a class redefinition error.
\begin{minted}[autogobble]{cpp}
    #ifndef _ACCOUNT_H_ // if this is not defined
    #define _ACCOUNT_H_ // define it. Otherwise skip to #endif

    #endif
\end{minted}
Some compilers allow you to include-guard your files with a simple preprocessor directive called \mintinline{cpp}{#pragma}.
\begin{minted}[autogobble]{cpp}
    #pragma once 
    class Account {};

\end{minted}
Specify the class in a file ``Account.h'':
\begin{minted}[autogobble]{cpp}
    #ifndef _ACCOUNT_H_
    #define _ACCOUNT_H_

    class Account {
        private:
        double balance;
        public:
        void set_balance(double bal);
        double get_balance();
    };

    #endif
\end{minted}
Implement the class in ``Account.cpp''
\begin{minted}[autogobble]{cpp}
    #include "Account.h"
    void Account::set_balance(double bal) {
        balance = bal;
    }
    double Account::get_balance() {
        return balance;
    }
\end{minted}

\subsection{Example}
\begin{minted}[autogobble]{cpp}
    #include <iostream>
    using namespace std;

    class Account {
    private:
        string name;
        double balance;
    public:
        void set_balance(double bal) {balance = bal;}
        double get_balance() {return balance;}

        void set_name(string n);
        string get_name();
        bool deposit(double amount);
        bool withdraw(double amount);
    };

    void Account::set_name(string n) {
        name = n;
    }
    string Account::get_name() {
        return name;
    }
    bool Account::deposit(double amount) {
        balance += amount;
        return true;
    }
    bool Account::withdraw(double amount) {
        if (balance-amount >= 0) {
            balance -= amount;
            return true;
        }
        else {
            return false;
        }
    }


    int main() {
        Account frank_account;
        frank_account.set_name("Frank's account");
        frank_account.set_balance(1000.00);
        if (frank_account.deposit(200.00)) {
            cout << "Deposit OK" << endl;
        }
        else {
            cout << "Deposit Not allowed." << endl;
        }

        if (frank_account.withdraw(500.00)) {
            cout << "Withdraw OK" << endl;
        } 
        else {
            cout << "Not sufficient funds." << endl;
        }

        if (frank_account.withdraw(1500.00)) {
            cout << "Withdraw OK" << endl;
        }
        else {
            cout << "Not sufficient funds" << endl;
        }
        return 0;
    }
    /* OUTPUT:
    Deposit OK
    Withdraw OK
    Not sufficient funds
    */
\end{minted}



%----------------------------------------------------------------------------------------
\section{Constructors and destructors}
\begin{itemize}
    \item Constructors are special member methods that are invoked during object creation, they are commonly used for initialization, constructors have the same name as the class and can be overloaded. 
    
    \item Destructors are also special member methods, they bear the same name as the class proceeded with a tilde \verb|~|, destructors are invoked automatically when an object is destroyed, they have no return type and no parameters, only 1 destructor is alowed per class and these cannot be overloaded. They are useful to release memory, close files, and other resources. You can call the destructor using \mintinline{cpp}{delete}, however if you don't call it, once the object goes out of scope it will automatically be deleted by C++.
    \item Example of constructors and destructors:
    \begin{minted}[autogobble]{cpp}
        class Player {
            private:
                std::string name;
                int health;
                int xp;
            public:
                // Constructors.
                Player();
                Player(std::string name);
                Player(std::string name, int health, int xp);
                // Destructors.
                ~Player();
        };

        class Account {
            private:
                std::string name;
                double balance;
            public: 
                // Constructors.
                Account();
                Account(std::string name, double balance);
                Account(std::string name);
                Account(double balance);
                // Destructors.
                ~Account();
        };
    \end{minted}
\end{itemize}

\subsection{Creating objects}
\begin{itemize}
    \item Objects can be called with any type of initialization and depending on the types of the parameters C++ will determine which one to use.
        \begin{minted}[autogobble]{cpp}
            #include <iostream>
            using namespace std;
            class Player {
                private:
                    std::string name;
                    int health;
                    int xp;
                public:
                // Constructors
                Player() {
                    cout << "No args constructor called." << endl;
                }
                Player(std::string name) {
                    cout << "String arg constructor called." << endl;
                }
                Player(std::string name, int health, int xp) {
                    cout << "Three args constructor called" << endl;
                }
                // Destructor.
                ~Player() {
                    cout << "Destructor called for " << name << endl;
                }
                void set_name(std::string name) {
                    this->name = name;
                }
            };

            class Account {
                private:
                    std::string name;
                    double balance;
                public: 
                    // Constructors.
                    Account();
                    Account(std::string name, double balance);
                    Account(std::string name);
                    Account(double balance);
            };

            int main() {
                {
                    Player slayer;
                    slayer.set_name("Slayer");
                }
                /* OUTPUT:
                No args constructor called.
                Destructor called for Slayer
                */
                {
                    Player frank;
                    frank.set_name("Frank");
                    Player hero("Hero");
                    hero.set_name("Hero");
                    Player villain("Villain", 100, 12);
                    villain.set_name("Villain");
                }
                /* OUTPUT: Notice the objects are destroyed in reverse order because they are in a stack.
                No args constructor called.
                String arg constructor called.
                Three args constructor called
                Destructor called for Villain
                Destructor called for Hero
                Destructor called for Frank
                */
                {
                    Player *enemy = new Player; // No args constructor.
                    enemy->set_name("Enemy");
                    Player *level_boss = new Player("Level Boss", 1000, 300); // 3 args constructor called.
                    level_boss->set_name("Level Boss");
                    delete enemy; // destructor called.
                    delete level_boss; // destructor called.
                }
                /* OUTPUT:
                No args constructor called.
                Three args constructor called
                Destructor called for Enemy
                Destructor called for Level Boss
                */
                return 0;
            }
        \end{minted}
\end{itemize}


%----------------------------------------------------------------------------------------
\section{The default constructor}
\begin{itemize}
    \item A default constructor does not expect any arguments, it is also called a no-args constructor. If you write no constructor at all for the class then C++ will generate a default constructor that does nothing. This default constructor is the one that is called when you create or instantiate a new object with no arguments.
    \item We are always free to provide our no args constructor and in fact it is best practice to do so.
    \item When there is at least one constructor in the class, C++ will not generate a dummy constructor, if we still need the no args constructor we must define it ourselves.
\end{itemize}

\subsection{Example}
\begin{minted}[autogobble]{cpp}
    #include <iostream>
    using namespace std;

    class Player {
        // C++ will generate a dummy constructor behind the scenes since there is none defined.
        private:
            std::string name;
            int health;
            int xp;
        public:
        void set_name(std::string name_val) {
            name = name_val;
        }
        std::string get_name() {
            return name;
        };
    };

    class Player_ {
        // C++ will not generate a dummy constructor, we have provided at least one.
        private:
            std::string name;
            int health;
            int xp;
        public:
        Player_() { // this is the default no args initializer.
            name = "None";
            health = 100;
            xp = 3;
        }
        void set_name(std::string name_val) {
            name = name_val;
        }
        std::string get_name() {
            return name;
        };
    };

    int main() {
        Player frank;
        frank.set_name("Frank");
        cout << frank.get_name() << endl;
        Player_ frank_with_constructor;
        cout << frank_with_constructor.get_name() << endl;
        return 0;
    }
    /* OUTPUT:
    Frank
    None
    */
\end{minted}


%----------------------------------------------------------------------------------------
\section{Overloading constructor}
\begin{itemize}
    \item Classes can have as many constructors as necessary, each must have a unique signature. The compiler has to be able to determine which to call based on the initialization information provided when creating the objects, if there is any ambiguity, the compiler won't guess, it will instead generate a compiler error. Default constructors are no longer provided for you classes once you've declared at least one.
\end{itemize}

\subsection{Example}
\begin{minted}[autogobble]{cpp}
    #include <iostream>
    using namespace std;

    class Player {
        private:
        std::string name;
        int health, xp;
        public:
        Player();
        Player(std::string name_val);
        Player(std::string name_val, int health_val, int xp_val);
    };

    Player::Player() {
        name = "None";
        health = 0;
        xp = 0;
    }

    Player::Player(std::string name_val) {
        name = name_val;
        health = 0;
        xp = 0;
    }

    Player::Player(std::string name_val, int health_val, int xp_val) {
        name = name_val;
        health = health_val;
        xp = xp_val;
    }

    int main() {
        Player empty; // None, 0, 0
        Player hero {"Hero"}; // Hero, 0, 0
        Player villain {"Villain"}; // Villain, 0, 0
        Player frank {"Frank", 100, 4}; // Frank, 100, 4
        Player *enemy = new Player("Enemy", 1000, 0); // Enemy, 1000, 0
        delete enemy;
        return 0;
    }
\end{minted}


%----------------------------------------------------------------------------------------
\section{Constructor Initialization lists} 
\begin{itemize}
    \item So far, all data memeber values have been set in the constructor body. What we really want to do is have the member data values initialized to our values before the constructor body executes, this is more efficient.
    \item Constructor initialization lists are: 
        \begin{itemize}
            \item More efficient.
            \item Initialization list immediately follows the parameter list.
            \item Initializes the data members as the objects created.
            \item Order of initialization is the order of declaration in the class.
        \end{itemize}
    
    \item Take the following example that does NOT use constructor initialization lists:
        \begin{minted}[autogobble]{cpp}
            Player::Player() {
                name = "None";
                health = 0;
                xp = 0;
            }
        \end{minted}
        \begin{itemize}
            \item While this works it takes more time and is inefficient. A better way is to provide a constructor initialization list.
            \item By the time we get to the statement \mintinline{cpp}{name = "None";} the string object has already been constructed and initialized to an empty string, in other words \mintinline{cpp}{name = "None";} just assigns to an already initialized object and doesn't initialize the object.
            \item So in order to initialize the members as they are created as to have a true initialization we must use a constructor initialization list.
        \end{itemize}
    
    \item Take the following example that does use constructor initialization lists:
        \begin{minted}[autogobble]{cpp}
            Player::Player() : name{"None"}, health{0}, xp{0} {
            }
        \end{minted}
        \begin{itemize}
            \item The above actually initializes before doing anything, this initializes everything before the body of the constructor is ever excecuted.
            \item The data memebers will not be initialized exactly in the order provided in the constructor initialization list, they will be initialized in the order in which they are provided in the class declaration.
        \end{itemize}
\end{itemize}

\subsection{Changing previous initialization to use constructor initialization lists}
\begin{minted}[autogobble]{cpp}
    #include <iostream>
    using namespace std;

    class Player {
        private:
        std::string name;
        int health, xp;
        public:
        Player();
        Player(std::string name_val);
        Player(std::string name_val, int health_val, int xp_val);
    };

    Player::Player() : name{"None"}, health{0}, xp{0} {
    }

    Player::Player(std::string name_val) : name{name_val}, health{0}, xp{0} {
    }

    Player::Player(std::string name_val, int health_val, int xp_val) : name{name_val}, health{health_val}, xp{xp_val}  {
    }

    int main() {
        Player empty; // None, 0, 0
        Player hero {"Hero"}; // Hero, 0, 0
        Player villain {"Villain"}; // Villain, 0, 0
        Player frank {"Frank", 100, 4}; // Frank, 100, 4
        Player *enemy = new Player("Enemy", 1000, 0); // Enemy, 1000, 0
        delete enemy;
        return 0;
    }
\end{minted}


%----------------------------------------------------------------------------------------
\section{Delegating constructors}
\begin{itemize}
    \item Often the code for constructors is very similar, duplicated code can lead to errors. C++ allows delegating constructors which is code for one constructor can call another in the initialization list, this in turn avoids duplicating code.
    \item This only works it constructors are called using constructor initialization lists, you cannot call constructors explicitly from the body of another. 
    \item When you call other constructors, the body of the called constructors will excecute and then excecution will continue from the constructor.
    \item So remember, if you delegate the constructors, both bodies of the constructors will be excecuted, the called constructor's body will be excecuted first and then the rest of the constructor body is excecuted.
\end{itemize}

\subsection{Example}
\begin{minted}[autogobble]{cpp}
    #include <iostream>
    using namespace std;

    class Player {
        private:
        std::string name;
        int health, xp;
        public:
        Player();
        Player(std::string name_val);
        Player(std::string name_val, int health_val, int xp_val);
    };

    Player::Player() : Player {"None", 0, 0 } {
        cout << "No args constructor" << endl;
    }

    Player::Player(std::string name_val) : Player{name_val, 0, 0} {
        cout << "One arg constructor" << endl;
    }

    Player::Player(std::string name_val, int health_val, int xp_val) : name{name_val}, health{health_val}, xp{xp_val}  {
        cout << "Three args constructor" << endl;
    }

    int main() {
        Player empty; 
        Player hero {"Hero"}; 
        Player frank {"Villain", 100, 55}; 
        return 0;
    }
    /* OUTPUT:
    Three args constructor
    No args constructor
    Three args constructor
    One arg constructor
    Three args constructor
    */
\end{minted}


%----------------------------------------------------------------------------------------
\section{Constructor parameters and default values}
\begin{itemize}
    \item We can use default values to simplify our code and reduce the number of overloaded constructors.
    \item Same rules apply as we learned with non-member functions.
    \item When doing this it is important to provide unambiguous constructors, for example if we declare one constructor with all default values and then declare a no arg constructor we would get a compiler error because C++ would not know which one to call.
\end{itemize}

\subsection{Example}
\begin{minted}[autogobble]{cpp}
    #include <iostream>
    using namespace std;

    class Player {
        private:
        std::string name;
        int health, xp;
        public:
        Player(std::string name_val = "None", int health_val = 0, int xp_val = 0);
    };

    Player::Player(std::string name_val, int health_val, int xp_val) : name{name_val}, health{health_val}, xp{xp_val}  {
        cout << "Three args constructor" << endl;
    }

    int main() {
        Player empty; 
        Player frank {"Frank"};
        Player hero {"Hero", 100};
        Player villain {"Villain", 100, 55};
        return 0;
    }
    /* OUTPUT:
    Three args constructor
    Three args constructor
    Three args constructor
    Three args constructor
    */
\end{minted}


%----------------------------------------------------------------------------------------
\section{Copy constructor}
\begin{itemize}
    \item When objects are copied C++ must create a new object from an existing object.
    \item When is a copy of an object made?
        \begin{itemize}
            \item Passing object by value as parameter.
            \item Returning an object from a function by value.
            \item Constructing one object based on another of the same class.
        \end{itemize}
    
    \item C++ must have a way of accomplish this copying, so C++ uses a copy constructor. C++ will provide a compiler defined copy constructor if you don't.
\end{itemize}

\subsection{1st usecase: Pass object by-value}
Example:
\begin{minted}[autogobble]{cpp}
    Player hero {"Hero", 100, 20};
    void display_player(Player p) {
        // p is a copy of hero in this example.
        // use p
        // Destructor for p will be called.
    }
    display_player(hero);
\end{minted}

\subsection{2nd usecase: Return object by value}
Example:
\begin{minted}[autogobble]{cpp}
    Player enemy;
    Player create_super_enemy() {
        Player an_enemy {"Super Enemy", 1000, 1000};
        return an_enemy; // A copy of an_enemy is returned, the copy is made by the copy constructor.
    }
    enemy = create_super_enemy();
\end{minted}

\subsection{Construct one object based on another}
Example:
\begin{minted}[autogobble]{cpp}
    Player hero{"Hero", 100, 100};
    Player another_hero{hero}; // A copy of hero is made.
\end{minted}

\subsection{Copy constructors}
\begin{itemize}
    \item If you don't provide your own way of copying objects by value then the compiler provides a default way of copying objects.
    \item The copy constructors copy the values of ach data memeber to the new object, this is called 'default-memberwise-copy' this basically means that for each data member C++ will copy them to the new object.
    \item In many cases this is fine, except for one thing, pointers.
    \item Beware if you have a pointer data member, because the pointer will be copied but data that the pointer is pointing to will not. This is explained in the next video.
\end{itemize}

\subsection{Best practice is:}
\begin{itemize}
    \item Provide a copy constructor when your class has raw pointer members.
    \item Provide the copy constructor with a const reference parameter.
    \item Use STL classes as they already provide copy constructors.
    \item Avoid using raw pointer data members if possible.
\end{itemize}

\subsection{Declaring the copy constructor}
\begin{itemize}
    \item Copy constructors are declared like this:
        \begin{minted}[autogobble]{cpp}
            Type::Type(const type &source);
            Player::Player(const Player &source);
            Account::Account(const Account &source);
        \end{minted}
    
    \item Actually implementing the copy constructor is like this:
        \begin{minted}[autogobble]{cpp}
                Player::Player(const Player &source) : name{source.name}, health{source.health}, xp{source.xp} {
                }
                Account::Account(const Account &source) : name{source.name}, balance{source.balance} {
                }
        \end{minted}
\end{itemize}

\subsection{Example}
\begin{minted}[autogobble]{cpp}
    #include <iostream>
    using namespace std;

    class Player {
        private:
        std::string name;
        int health;
        int xp;
        public: 
        std::string get_name() {return name;}
        int get_health() {return health;}
        int get_xp() {return xp;}
        Player(std::string name_val = "None", int health_val = 0, int xp_val = 0);
        // Copy constructor:
        Player(const Player &source);
        // Destructor:
        ~Player() { cout << "Destrocutor called for: " << name << endl; }
    };

    Player::Player(std::string name_val, int health_val, int xp_val) : name{name_val}, health{health_val}, xp{xp_val} {
        cout << "Theree-args constructor for " + name << endl;
    }

    Player::Player(const Player &source) : name{source.name}, health{source.health}, xp{source.xp} {
        cout << "Copy constructor - made a copy of: " << source.name << endl;
    }

    void display_player(Player p) {
        cout << "Name: " << p.get_name() << endl;
        cout << "Health: " << p.get_health() << endl;
        cout << "xp: " << p.get_xp() << endl;
    }

    int main() {
        Player empty;
        display_player(empty);
        Player copy_of_empty {empty};
        display_player(copy_of_empty);
        Player frank {"Frank"};
        Player hero {"Hero", 100};
        Player villain {"Villain", 100, 55};
        return 0;
    }
    /* OUTPUT:
    Theree-args constructor for None
    Copy constructor - made a copy of: None
    Name: None
    Health: 0
    xp: 0
    Destrocutor called for: None
    Copy constructor - made a copy of: None
    Copy constructor - made a copy of: None
    Name: None
    Health: 0
    xp: 0
    Destrocutor called for: None
    Theree-args constructor for Frank
    Theree-args constructor for Hero
    Theree-args constructor for Villain
    Destrocutor called for: Villain
    Destrocutor called for: Hero
    Destrocutor called for: Frank
    Destrocutor called for: None
    Destrocutor called for: None
    */
\end{minted}

\section{Shallow Copying with the Copy Constructor}
\begin{itemize}
    \item Consider a class that contains a pointer as a data member.
    \item Constructor allocates storage dynamically and initializes the pointer, the destructor releases the memory allocated by the constructor.
    \item The default copy constructor conducts a memberwise copy of the data members, thus each data member is copied from the source object and the pointer's value is copied, not the values it points to, thus we have two objects that point to the same values and the change in the pointer's values can be done using both objects. This can lead to all sorts of problems for example if one object's destructor is called and releases the memory the other object will still point to the released storage.
\end{itemize}

\subsection{Example}
\begin{minted}[autogobble]{cpp}
    #include <iostream>
    using namespace std;

    class Shallow {
        public:
            int *data;
        Shallow(int d); // Constructor.
        Shallow(const Shallow &Source); // Copy constructor.
        ~Shallow(); // Destructor.
    };

    Shallow::Shallow(int d) {
        data = new int; // allocate storage.
        *data = d;
    }

    Shallow::~Shallow() {
        delete data; // free storage.
        cout << "Destructor freeing data" << endl;
    }

    Shallow::Shallow(const Shallow &source) : data(source.data) {
        cout << "Copy constructor - Shallow" << endl;
    }

    void display_shallow(Shallow obj) { // obj is a copy of obj1.
        cout << *obj.data << endl;
        // obj goes out of scope and is thus destroyed. The pointer is freed.
    }

    int main() {
        Shallow obj1 {100};
        display_shallow(obj1);
        cout << *obj1.data << endl; // Referencing a pointer that has already been freed.

        return 0;
    }
    /* OUTPUT:
    Copy constructor - Shallow
    100
    Destructor freeing data   
    15869168 // trash data
    */
\end{minted}

\section{Deep Copying with the Copy constructor}
\begin{itemize}
    \item Create a copy of the pointed-to data. Each copy will have a pointer to unique storage in the heap. Deep copy when you have a raw pointer as a class member.
    \item You can also delegate constructors from copy constructors.
\end{itemize}

\subsection{Example}
\begin{minted}[autogobble]{cpp}
    #include <iostream>
    using namespace std;

    class Deep {
        public:
        int *data;
        Deep(int d);
        Deep(const Deep &source);
        ~Deep();
    };

    Deep::Deep(int d) {
        data = new int; // allocate storage.
        *data = d;
    }

    Deep::~Deep() {
        delete data; // free storage.
        cout << "Destructor freeing data" << endl;
    }

    // Deep::Deep(const Deep &source) {
    //     data = new int;
    //     *data = *source.data;
    //     cout << "Copy constructor -deep" << endl;
    // }

    Deep::Deep(const Deep &source) : Deep{*source.data} /* Delegating to the first constructor */ {
        cout << "Copy constructor -deep" << endl;
    }

    void display_deep(Deep s) {
        cout << *s.data << endl;
    }

    int main() {
        Deep obj1 {100};
        display_deep(obj1);
        *obj1.data = 100;
        return 0;
    }
    /* OUTPUT:
    Copy constructor -deep
    100
    Destructor freeing data
    Destructor freeing data
    */
\end{minted}


%----------------------------------------------------------------------------------------
\section{Move constructor}
\begin{itemize}
    \item Sometimes when we excecute the compiler creates unnamed temporary values:
        \begin{minted}[autogobble]{cpp}
            int total {0};
            total = 100 + 200;
        \end{minted}
    
    \item In the above, 100 + 200 is not addressable, it is a right-hand side value which then gets assigned to the left-hand side which is the variable 'total' and the temporary objects.
    \item Naturally, the problem here is that this temporary values can be very big and with objects, shallow copies and deep pointer copies the overhead can get really big and this is why we need the move constructors come in.
    \item Sometimes copy constructors are called many times automaticallt due to the copy semantics of C++.
    \item Copy constructors doing deep copying can have a significant performance (bottleneck).
    \item C++11 introduced move semantics and the move constructor, the move constructor moves an object rather than copying it.
    \item Optional but recomended when you have raw pointers.
    \item Copy elision : C++ may optimize copying away completely (RVO- Return Value Optimization).
    \item If you don't provide a move construtor then what will be used instead is the copy constructor, for efficiency you should always provide a move constructor.
\end{itemize}

\subsection{r-value references}
\begin{itemize}
    \item Used in moving semantics and perfect forwarding.
    \item Move semantics is all about r-value references.
    \item Used by move constructor and move assignment operator to efficiently move an objet rather than copy it. 
    \item R-value reference operator (\&\&).
    \item Example of l-value reference parameters:
        \begin{minted}[autogobble]{cpp}
            int x {10};
            int &l_ref = x; // l-value reference.
            l_ref = 10; // change x to 10.

            int &&r_ref = 200; // r-value ref.
            r_ref = 300; // change r_ref to 300.

            int &&x_ref = x; // compiler error.
        \end{minted}
    
    \item Another example:
        \begin{minted}[autogobble]{cpp}
            int x {100}; // x is an l-value.
            void func(int &num); // Define function that expects an l-value.
            func(x); // x is an l-value so OK.
            func(200); // 200 is a r-value.
        \end{minted}
    
    \item Example of r-value reference parameters:
        \begin{minted}[autogobble]{cpp}
            int x {100}; // x is an l-value.
            void func(int &&num); // A funcion that expects an r-value reference.
            func(200); // calls func with an r-value.
            func(x); // error: x is an l-value.
        \end{minted}
    
    \item We could, in any case, overload the functions.
        \begin{minted}[autogobble]{cpp}
            int x {100};
    
            void func(int &num); // function that expects a l-value.
            void func(int &&num); // function that expects a r-value.

            func(x); // calls: func(int &num);
            func(200); // calls: func(int &&num);
        \end{minted}
\end{itemize}

\subsection{Move constructors}
\begin{itemize}
    \item Move constructors basically is the memberwise copy of objects but the pointer that is copied in one of the objects is nulled out in the other, so just one pointer is being used, you just moved a pointer not all the deep copy things the pointers where pointing to.
    \item Another way of thinking about it is that the move constructor just steals the data and then nulls out the source pointer.
\end{itemize}
\begin{minted}[autogobble]{cpp}
    #include <iostream>
    #include <vector>
    using namespace std;

    class Move {
        public: 
        int *data;
        Move(int d); // Constructor.
        Move(const Move &source); // Copy constructor.
        ~Move(); // Destructor.
        Move(Move &&source) noexcept; // Move constructor.
    };

    Move::Move(int d) {
        data = new int;
        *data = d;
        cout << "Constructor called: " << *data << endl;
    }

    Move::Move(const Move &source) : Move(*source.data) {
        cout << "Copy constructor called -deep copy: " << *data << endl;
    }

    Move::Move(Move &&source) noexcept : data{source.data} {
        source.data = nullptr;
        cout << "Move constructor - moving resource: " << *data << endl;
    }

    Move::~Move() {
        cout << ((data == nullptr)? "Destructor called for nullptr.": "Destructo called.") << endl;
        delete data;
    }

    int main() {
        vector<Move> vec;
        // Move{10} is un-named, C++ will use the copy constructors to make copies, 
        // but since we have provided a move constructor C++ wont use the copy cosntructor.
        vec.push_back(Move{10}); 
        vec.push_back(Move{20}); 
        return 0;
    }
    /* OUTPUT:
    Constructor called: 10
    Move constructor - moving resource: 10
    Destructor called for nullptr.
    Constructor called: 20
    Move constructor - moving resource: 20
    Move constructor - moving resource: 10
    Destructor called for nullptr.
    Destructor called for nullptr.
    Destructo called.
    Destructo called.
    */
\end{minted}


%----------------------------------------------------------------------------------------
\section{The 'this' pointer}
\begin{itemize}
    \item 'this' is a reserved word that referes to the current object.
    \item Contains the address of the object, so it's a pointer to the object.
    \item Can only beused in class scope.
    \item All members acces is done via the 'this' pointer.
    \item Can be used bu the programmer to: 
        \begin{itemize}
            \item Access data members and methods.
            \item To determine if two objects are the same (later on).
            \item Can be dereferences (*this) to yield the current object.
        \end{itemize}
    
    \item This is also used to have same name parameters and class data members.
        \begin{minted}[autogobble]{cpp}
            void account::set_balance(double balance) {
                balance = balance; // which balance? The parameter or the data memeber?
            }

            void account::set_balance(double balance) {
                this->balance = balance; // Now it is unambigous.
            }
        \end{minted}
    
    \item To determine object identity (this will be useful to overload operators.):
        \begin{minted}[autogobble]{cpp}
            int Account::compare_balance(const Account &other) {
                if (this == &other) {
                    std::cout << "The same objects" << std::endl;
                }
            }
            frank_account.compare(frank_account);
        \end{minted}
\end{itemize}

%----------------------------------------------------------------------------------------
\section{Using const with classes}
\begin{itemize}
    \item const classes.
        \begin{itemize}
            \item Pass arguments to class member methods as \mintinline{cpp}{const}.
            \item We can also create const objects.
            \item What happens if we call member functions on const objects?
            \item const-correctness. 
        \end{itemize}
    
    \item Creating a const object:
        \begin{itemize}
            \item villain is a const object so it's attributes cannot change.
                \begin{minted}[autogobble]{cpp}
                    const Player villain {"Villain", 100, 55};
                \end{minted}
            
            \item What happens when we call member methods on a const object?
                \begin{minted}[autogobble]{cpp}
                    const Player villain {"Villain", 100, 55};
                    void display_player_name(const Player &p) {
                        cout << p.get_name() << endl;
                    }
                    display_player_name(villain); // Error. Even if we are not changing the variables and data members the compiler does not allow it.
                \end{minted}
            
            \item The compiler assumes that any method called can potentially modify the object, so in the class declaration we must declare methods as const as well, that way we tell the compiler that we are sure that that method will not change the object. Once you have declared a method to be const if you change any data member in the method body the compiler will produce an error.
                \begin{minted}[autogobble]{cpp}
                    class Player {
                        private:
                        ...
                        public: 
                        ...
                        std::string get_name() const; // This method will not modify the data members.
                    };
                    const Player villain {"Villain", 100, 55}; 
                    villain.set_name("Nice guy"); // Not allowed.
                    std::cout << villain.get_name() << std::endl; // OK
                \end{minted}
        \end{itemize}
\end{itemize}

%----------------------------------------------------------------------------------------
\section{Static Class Members}
\begin{itemize}
    \item Class data members can be declared as static.
        \begin{itemize}
            \item A single data member that belongs to the class, not the objects.
            \item Useful to store class-wide information.
            \item Example: a variable that counts the instances of the class, it cannot be a member variable it needs to be a static member.
        \end{itemize}
    \item Class functions can be declared as static.
        \begin{itemize}
            \item Independent of any objects.
            \item Can be called using the class name.
        \end{itemize}
    
    \item Example of counting the number of players instanciated in the program:
        \begin{minted}[autogobble]{cpp}
            class Player {
                private: 
                    static int num_players;
                public:
                    static int get_num_players() {
                        ...
                    }
            };
        \end{minted}
    
    \item A static variable is only intialized ONCE. This happens tipically in the cpp file not the header. You can only instantiate the static variables outside the class, not inside.
        \begin{minted}[autogobble]{cpp}
            #include "Player.h"
            int Player::num_players = 0; // initialization (only once)
        \end{minted}
    
    \item Static methods only have access to static data. It does not access to non-static data members.
        \begin{minted}[autogobble]{cpp}
            int Player::get_num_players() {
                return num_players;
            }
        \end{minted}
    
    \item Lets write an example of how to increment and decrement a static variable that counts the instances of an object.
        \begin{minted}[autogobble]{cpp}
            Player::Player(std::string name_val, int health_val, int xp_val) : name{name_val}, health{health_val}, xp{xp_val} {
                ++num_players;
            }
            Player::~Player() {
                --num_players;
            }
        \end{minted}
        \begin{itemize}
            \item Be careful to only have the increment in one constructor, since if you delegate constructors you might be incrementing/decrementing the static variable by more than one.
        \end{itemize}
    
    \item Now by calling the method ``display\_active\_players'' we can check how many instances we have of the Player class.
        \begin{minted}[autogobble]{cpp}
            void display_active_players() {
                cout << "Active players:" << Player::get_num_players() << endl;
            }
        \end{minted}
\end{itemize}

\section{Example}
In ``Player.h''
\begin{minted}[autogobble]{cpp}
    #ifndef PLAYER_H 
    #define PLAYER_H
    #include <string>
    class Player {
        private:
            static int num_players;
            std::string name;
            int health;
            int xp;
        public: 

        std::string get_name() {return name;}
        int get_health() {return health;}
        int get_xp() {return xp;}
        
        Player(std::string name_val = "None", int health_val = 0, int xp_val = 0);
        // Copy constructors
        Player(const Player &source) {}
        // Destructor
        ~Player();

        static int get_num_players(); // does not have access to name, health, or xp
    }
\end{minted}
In ``Player.cpp''
\begin{minted}[autogobble]{cpp}
    #include "Player.h"
    int Player::num_players {0};

    Player::Player(std::string name_val, int health_val, int xp_val) 
        : name{name_val}, health{health_val}, xp{xp_val} {
        ++num_players;
    }
    Player::Player(const Player &source) {
    }
    Player::~Player() {
        --num_players;
    }
    int Player::get_num_players() {
        return num_players;
    }
\end{minted}
In ``main.cpp''
\begin{minted}[autogobble]{cpp}
    #include <iostream>
    #include "Player.h"
    using namespace std;

    void display_active_players() {
        cout << "Active players: " << Player::get_num_players() << endl;
    }

    int main() {
        display_active_players();
        Player hero{"Hero"};
        display_active_players();
        {
            Player frank{"Frank"};
            display_active_players();
        }
        display_active_players();
        return 0;
    }
    /* OUTPUT:
    0
    1
    2
    1
    */
\end{minted}

%----------------------------------------------------------------------------------------
\section{Structs vs Classes}
\begin{itemize}
    \item \mintinline{cpp}{struct} is supported in C++ and it comes from C.
    \item In addition to define a \mintinline{cpp}{class} we can declare a \mintinline{cpp}{struct}.
    \item Essentially the same as a \mintinline{cpp}{class} except that all the members are public by default.
\end{itemize}

\subsection{When to use a class or struct?}
\begin{itemize}
    \item \mintinline{cpp}{struct}
        \begin{itemize}
            \item Use \mintinline{cpp}{struct} for passive objects with public access.
        \end{itemize}
    \item \mintinline{cpp}{class}
        \begin{itemize}
            \item Use class for active objects with passive access.
            \item Implementing getters/setters as needed.
            \item Implementing member methods as needed.
        \end{itemize}
\end{itemize}

\subsection{Example}
\begin{itemize}
    \item \mintinline{cpp}{class}
        \begin{minted}[autogobble]{cpp}
            struct Person {
                std::string name; // public
                std::string get_name(); // public 
            }
            Person p;
            p.name = "Frank"; // OK
            std::cout << p.get_name(); // OK
        \end{minted}
        
    \item \mintinline{cpp}{struct}
        \begin{minted}[autogobble]{cpp}
            class Person {
                std::string name; // private
                std::string get_name(); // private
            };
            Person p;
            p.name = "Frank"; // compiler error, data is private.
            std::cout << p.get_name(); // compiler error, data is private.
        \end{minted}
\end{itemize}


%----------------------------------------------------------------------------------------
\section{Friends of a class}
\begin{itemize}
    \item Friends are a very big controversial in C++, the whole argument is wether friends aid encasulation or enhance it. When a class extends friendship to a fucntion of a class then the other class (the one being granted friendship) can now access the private class members.
    \item Friend:
        \begin{itemize}
            \item A function of class that has access to private class members.
            \item And, that function of or class is NOT a member of the class it is accessing.
        \end{itemize}
    
    \item Function:
        \begin{itemize}
            \item Can be regular non-member functions.
            \item Can be member methods of another class.
        \end{itemize}
    
    \item Class:
        \begin{itemize}
            \item Another class can have access to private class members.
        \end{itemize}
\end{itemize}

\subsection{Friends of a class}
\begin{itemize}
    \item Friendship must be granted not taken:
        \begin{itemize}
            \item Declared explicitly in the class that is granting friendship.
            \item Declared in the function prototype with the keyword friend.
        \end{itemize}
    
    \item Friendship is not symmetric. (not reciprocal)
        \begin{itemize}
            \item Must be explicitly granted.
                \begin{itemize}
                    \item If A is a friend of B.
                    \item B is NOT a friend of A.
                \end{itemize}
        \end{itemize}
    
    \item Friendship is not transitive:
        \begin{itemize}
            \item Must be explicitly granted.
                \begin{itemize}
                    \item If A is a friend of B AND 
                    \item B is a friend of C 
                    \item Then A is NOT a friend of C.
                \end{itemize}
        \end{itemize}
    
    \item Friendship is not inherited.
    
    \item Example of non-member friend functions:
        \begin{minted}[autogobble]{cpp}
            class Player {
                friend void display_player(Player &p);
                std::string name;
                int health;
                int xp;
                public:
                ...
            };
            // Now lets declare a non-member function that will change the private members of Player.
            void display_player(Player &p) {
                std::cout << p.name << std::endl;
                std::cout << p.health << std::endl;
                std::cout << p.xp << std::endl;
            }
        \end{minted}
    
    \item Example of other class member function friends:
        \begin{minted}[autogobble]{cpp}
            class other_class {
                ...
                public:
                void display_player(Player &p) {
                    std::cout << p.name << std::endl;
                    std::cout << p.health << std::endl;
                    std::cout << p.xp << std::endl;
                }
            };
        \end{minted}
    
    \item Declaring an entire class a friend.
        \begin{minted}[autogobble]{cpp}
            class Player {
                friend class other_class;
                std::string name;
                int health;
                int xp;
            };
        \end{minted}
    
    \item Friendship allows you to not write setters and getters just for using the attributes in one other class, it allows you to grant other classes permision to change those variables without having to write setters and getters.
    \item Friendship doesn't always only serve to save time with writing setters/getters, it works for all sorts of things and must only be used when they make sense.
    \item Only the minimal friendship should be granted.
\end{itemize}
